% ============================================================
% Capitulo 3 -- Metodologia
% Tamanho-alvo: 8-12 paginas
% Ordem de escrita: PRIMEIRO (junto com Resultados)
% ============================================================

\chapter{Metodologia}
\label{cap:metodologia}

% ============================================================
\section{Vis\~ao geral do pipeline}
\label{sec:pipeline}

% TESE: pipeline em 5 blocos (dados brutos -> grafos -> features ->
%       modelo -> avaliacao), mesmo para os 3 datasets.
%
% ANTES: 1 paragrafo descrevendo os 5 blocos do diagrama.

\begin{figure}[H]\centering
\resizebox{\linewidth}{!}{%
\begin{tikzpicture}[
    node distance=0.9cm and 0.7cm,
    box/.style    ={rectangle, rounded corners, draw, thick,
                    minimum width=2.0cm, minimum height=0.9cm,
                    align=center, font=\scriptsize},
    data/.style   ={box, fill=blue!10},
    proc/.style   ={box, fill=green!15},
    model/.style  ={box, fill=orange!18},
    eval/.style   ={box, fill=red!12},
    arr/.style    ={->, >=stealth, thick}
]
\node[data]  (raw)    {Dados brutos\\FNN, UPFD, Bluesky};
\node[proc, right=of raw]   (graph)  {Constru\c{c}\~ao\\de grafos};
\node[proc, right=of graph] (feat)   {\emph{Features}\\posicionais\\+\,BERT};
\node[model, right=of feat] (model)  {GNN\\GCN/GAT/SAGE};
\node[eval, right=of model] (val)    {Valida\c{c}\~ao\\k-fold pareado};
\node[eval, right=of val]   (out)    {M\'etricas\\F1, Cohen's $d$};

\draw[arr] (raw)   -- (graph);
\draw[arr] (graph) -- (feat);
\draw[arr] (feat)  -- (model);
\draw[arr] (model) -- (val);
\draw[arr] (val)   -- (out);
\end{tikzpicture}}
\caption{Vis\~ao geral do \emph{pipeline} de detec\c{c}\~ao.}
\label{fig:pipeline}
\end{figure}

% DEPOIS: nao precisa interpretar -- a figura fala por si.
% TRANSICAO: anunciar que as proximas secoes detalham cada bloco.

% ============================================================
\section{Datasets}
\label{sec:datasets}

% TESE: 3 datasets, papeis distintos -- FNN/UPFD para validacao
%       supervisionada, Bluesky para aplicacao sem labels.
%
% ANTES: 1 paragrafo de overview.

\begin{table}[H]\centering
\caption{Datasets utilizados, com tamanho, fonte e papel no trabalho.}
\label{tab:datasets}
\small
\begin{tabular}{lrrll}
\toprule
\textbf{Dataset} & \textbf{N grafos} & \textbf{N classes} & \textbf{Fonte} & \textbf{Papel} \\
\midrule
FakeNewsNet PolitiFact   & $\sim$754   & 2 (50/50)   & \cite{shu2020fakenewsnet}      & Treino + valida\c{c}\~ao \\
UPFD-PolitiFact          & $\sim$314   & 2 (50/50)   & \cite{dou2021upfd}             & Cross-dataset \\
UPFD-GossipCop           & $\sim$5{.}464 & 2 (50/50) & \cite{dou2021upfd}             & Cross-dataset \\
Bluesky (acad\^emico)    & 168{.}463   & --          & \cite{failla2024bluesky}       & Aplica\c{c}\~ao (sem \emph{labels}) \\
\bottomrule
\end{tabular}
\end{table}

% DEPOIS: 1 paragrafo por dataset (4 paragrafos curtos).
% - FNN: coletado por Shu, cascatas estrela por falta de timestamps.
% - UPFD-PolitiFact: 314 grafos, balanceado.
% - UPFD-GossipCop: 5464 grafos, dominio celebridade/entretenimento.
% - Bluesky: 168k posts, sem ground truth -> modo demonstracao no Cap 5.

% ============================================================
\section{Constru\c{c}\~ao de grafos: FakeNewsNet}
\label{sec:construcao_fnn}

% TESE: 00_construir_grafos_fakenewsnet.py constroi estrela com 3 variantes
%       de features posicionais (posfull, posmin, posgrau).
%
% CONTEUDO (~1.5 pagina):
% - Pipeline: CSV -> BERT do titulo -> grafo (raiz=noticia, filhos=tweets) ->
%   features posicionais -> cache em data/fakenewsnet_posfull/.
% - 3 variantes: posfull (768+3), posmin (768+1), posgrau (768+2).
% - Justificar estrela: ausencia de parent_of_retweet no CSV bruto.

% ============================================================
\section{Folds compartilhados}
\label{sec:folds}

% TESE: 10 folds estratificados gerados UMA vez (seed=42), reusados em
%       todos os experimentos -- pre-requisito do ttest_rel.
%
% CONTEUDO (~0.5 pagina):
% - gerar_folds.py com StratifiedKFold(n_splits=10, random_state=42).
% - Salvos em data/folds_fnn.pt.
% - Justificar: modelos comparados precisam estar nos mesmos folds.

% ============================================================
\section{Arquiteturas GNN e hiperpar\^ametros}
\label{sec:arquiteturas}

% TESE: 3 arquiteturas com hiperparametros identicos -- isola o efeito
%       arquitetural do efeito de tuning.
%
% ANTES: 1 paragrafo introducindo a tabela.

\begin{table}[H]\centering
\caption{Hiperpar\^ametros usados nos modelos GNN.}
\label{tab:hyperparams}
\small
\begin{tabular}{lrrr}
\toprule
\textbf{Hiperpar\^ametro} & \textbf{GCN} & \textbf{GAT} & \textbf{SAGE} \\
\midrule
\texttt{hidden\_channels}     & 64     & 64     & 64 \\
\texttt{num\_layers}          & 3      & 3      & 3 \\
\texttt{dropout}              & 0{.}5  & 0{.}5  & 0{.}5 \\
\texttt{num\_heads} (atten\c{c}\~ao) & --     & 4      & -- \\
\texttt{pooling}              & \texttt{global\_mean\_pool} & idem & idem \\
\texttt{learning\_rate}       & 0{.}001 & 0{.}001 & 0{.}001 \\
\texttt{weight\_decay}        & $5\times10^{-4}$ & $5\times10^{-4}$ & $5\times10^{-4}$ \\
\texttt{batch\_size}          & 32     & 32     & 32 \\
\texttt{optimizer}            & Adam   & Adam   & Adam \\
\texttt{scheduler}            & \multicolumn{3}{c}{ReduceLROnPlateau (factor=0{.}5, patience=5)} \\
\texttt{early\_stopping}      & \multicolumn{3}{c}{patience=7 sobre F1-macro de valida\c{c}\~ao} \\
\texttt{max\_epochs}          & 30     & 30     & 30 \\
\texttt{seed}                 & 42 (com varredura para 10 \emph{seeds} em \S\ref{sec:topo_sem_texto}) \\
\bottomrule
\end{tabular}
\end{table}

% DEPOIS: 1 paragrafo justificando os defaults.

% ============================================================
\section{\emph{Baselines} n\~ao-GNN}
\label{sec:baselines}

% TESE: 2 baselines -- LogReg-BERT (teto textual) e RF tabular
%       (teto inferior estrutural).
%
% CONTEUDO (~1 pagina):
% - LogReg-BERT-root: input = x[0] do grafo (BERT do titulo da raiz, 768 dim).
% - RF-tabular: input = [num_nodes, grau_root] (2 numeros agregados).
% - Justificativa: cada um isola uma direcao do espaco de hipoteses.

% ============================================================
\section{Protocolo de avalia\c{c}\~ao}
\label{sec:protocolo}

% TESE: F1-macro como metrica primaria; k-fold pareado + ttest_rel para
%       comparacoes de modelo; Cohen's d para analise estrutural.
%
% CONTEUDO (~1 pagina):
% - F1-macro vs accuracy: justificar (classes balanceadas, evita inflacao).
% - ttest_rel pareado: usado em §4.2 e §4.3.
% - Cohen's d: usado em §4.5.
% - Reprodutibilidade caveat: CPU-only, seeds fixadas, nao bit-exact.

% ============================================================
\section{An\'alise estratificada do erro}
\label{sec:metod_estratificacao}

% TESE: 3 tecnicas pra cruzar metrica estrutural com acerto do modelo --
%       LOWESS (continua), percentil empirico (categorica), bootstrap (CI).
%
% CONTEUDO (~0.7 pagina):
% - LOWESS via rolling-mean (janela 10% da amostra, minimo 30): visualiza
%   P(acerto | metrica) sem cherry-picking de bins.
% - Percentil empirico (vs z-score): distribuicoes sao heavy-tail ->
%   outlier via sigma e' mal-definido.
% - Bootstrap pareado (n_boot=2000): CI 95% para diff F1(SAGE)-F1(RF) em
%   subsets nao-pareados entre folds.
%
% CITACOES: \cite{cleveland1979lowess}, \cite{cohen1988statistical},
%           \cite{efron1993bootstrap}.

% ============================================================
\section{An\'alise de explicabilidade}
\label{sec:analise_explicabilidade}

% TESE: GNNExplainer aplicado em SAGE estrutural (UPFD-GossipCop),
%       20 amostras estratificadas por (classe x tamanho).
%
% CONTEUDO (~0.5 pagina):
% - Setup: torch_geometric.explain.Explainer, epochs=200, edge_mask.
% - Analise por hop-distance (1, 2, 3+).
% - Estratificacao: 5 fake-pequena, 5 fake-grande, 5 real-pequena, 5 real-grande.

% ============================================================
\section{Classificador dual: regra de combina\c{c}\~ao \emph{post-hoc}}
\label{sec:dual}

% TESE CRITICA: regra dual e' heuristica POST-HOC, derivada APOS
%       observar resultados de §4.6. Nao e' decisao a priori.
%
% CONTEUDO (~0.5 pagina):
% - Anunciar que a aplicacao do Cap 5 combina textual + topologico
%   via regra empirica.
% - Justificativa empirica e detalhes ficam em \S\ref{sec:concordancia}.
% - NAO escrever a equacao aqui -- so anunciar.

% ============================================================
\section{Reprodutibilidade}
\label{sec:reprodutibilidade}

% TESE: requirements pinado + seed=42 universal -> reproduzivel exceto
%       por bit-exactness (CUDA/cuDNN).
%
% CONTEUDO (~0.5 pagina):
% - requirements.txt pinado.
% - Python 3.12.10, torch 2.10.0+cpu.
% - Seeds: 42 universal; 10 seeds para variancia em \S\ref{sec:topo_sem_texto}.
% - Citar README.md como fonte autoritativa.
